\section{Consequences and remaining limits}\label{sec:conclusion}
Affine separator certificates can be constructed without a known
optimizer by alternating a lower-bound calculation and a new graded
partition. The decisive estimate controls the sum of squared disagreements
between coordinate copies. It gives a contraction on branching trees
while preserving a final partition count that is logarithmic in accuracy.
Certified local errors can be incorporated without charging an error for
every tested pair of boxes. For polynomial objectives, affine Taylor
models turn the abstract construction into a rational algorithm with
explicit local checks.

The nonlinear dynamics extension shows how a coupling constraint changes
the construction. Forward repair controls feasibility, but it must be
paired with adjoint-adjusted slopes to cancel first-order state errors.
Contractivity bounds the repair and adjoint constants uniformly in the
horizon. The finite-precision analysis separates an exactly feasible
recurrence representation from the potentially much larger expansion of
all rational states.

These results concern certified accuracy under stated growth and
conditioning assumptions. They give neither a general complexity bound
for all sparse nonconvex problems nor a performance comparison with
existing solvers. Several questions remain beyond the present theorems:
how to exploit several minimizers without imposing point growth, whether
weaker stability conditions suffice for nonlinear dynamics, and whether
pointwise projected covering lower bounds have matching constructive
upper bounds in more than one separator dimension. Resolving these
questions requires additional arguments; none is used in the guarantees
proved here.
